\section{Evaluation}\label{sec:evaluation}

\subsection{Study design and recorded versions}\label{sec:evaluation-design}

The evaluation asks whether checking feedback helps author usable templates, what template reuse costs, and how closely the educational assessments match their references. The latest authoring comparison is reported first. Earlier studies used different requests or implementation versions and remain supporting evidence in Appendix~\ref{app:comparisons}.

Each study records its model, provider, prompts and settings. Aggregation is limited to the configurations within a stated study; results from earlier runs are not added to the latest totals. Appendix~\ref{app:history} identifies the source revisions and saved records.

The main study used fifteen distinct requests: five Python concepts at three requested difficulty levels. Each request was evaluated with three OpenAI models, giving forty-five paired model--request comparisons. The three strategies shared the initial proposal. One-shot generation ended there. Check-and-regenerate could request fresh proposals, while feedback-guided revision received the previous candidate and check findings. Both checked strategies retained the initial plan and allowed three candidate opportunities in total, stopping on successful workflow completion. Malformed responses consumed an opportunity. The experimental wrapper could make a fresh attempt when no valid candidate was available; this differs from the production default of one initial proposal plus three revisions.

Both checked strategies used the application's technical checks and difficulty policy. Independent offline evaluation was performed after the branches finished, so its findings were never revision feedback. It used 64 generated examples, a maximum of 192 execution calls, a 30-second checking budget and seed 20261102. Local previews used seeds 7, 19 and 73. The primary outcome was whether each job's final available candidate passed these technical checks, including candidates withheld from delivery. Recorded request constraints and workflow completion were separate outcomes.

Authoring requested low reasoning, a 16,384-token completion allowance and a 300-second HTTP deadline. The common SoCLaaS difficulty judge used provider-default reasoning. Equal candidate allowances did not imply equal token usage or latency. Uncertainty estimates grouped the three model results for each request together, using fifteen request clusters and 20,000 bootstrap resamples. These results describe the tested panel, rather than a random sample of all models or authoring tasks.

\subsection{Technical validity, repair and delivery}\label{sec:technical-results}\label{sec:benefits}\label{sec:delivery}

All forty-five OpenAI comparisons completed. Table~\ref{tab:authoring} separates candidates passing offline checks from those that completed the authoring workflow. All technical passes also met the recorded machine-readable request constraints.

\begin{table}[!htbp]
\centering\setstretch{1}\small
\setlength{\tabcolsep}{5pt}\renewcommand{\arraystretch}{1.15}
\caption{Latest OpenAI authoring comparison: 45 jobs per strategy}\label{tab:authoring}
\begin{tabularx}{\linewidth}{@{}X P{.14\linewidth}P{.19\linewidth}P{.19\linewidth}@{}}
\toprule
\textbf{Outcome} & \textbf{One-shot} & \textbf{Regenerate} & \textbf{Revise}\\
\midrule
Final technical passes & 14 (31.1\%) & 22 (48.9\%) & 30 (66.7\%)\\
Delivered technical passes & --- & 21 & 29\\
Technical passes withheld & --- & 1 & 1\\
Author model calls & 45 & 101 & 100\\
Difficulty-judge calls & 0 & 25 & 33\\
Median strategy time (seconds) & 12.02 & 93.54 & 152.05\\
\bottomrule
\end{tabularx}
\end{table}

Feedback revision produced sixteen more technical passes than one-shot generation, a gain of 35.6 percentage points (request-cluster 95\% interval: 24.4 to 46.7). No initially passing candidate became a final technical failure in either checked strategy. The difference between revision and regeneration was smaller: eight additional passes, or 17.8 points (interval: 0 to 35.6; request-cluster swap test $p\approx0.117$). Revision performed better descriptively, but this study leaves its advantage over regeneration uncertain. The comparison also changes both access to the previous candidate and detailed failure feedback; it does not isolate feedback alone.

Distractor selection was the most frequent offline failure: twenty-three one-shot candidates and twelve final revision candidates failed to provide a common usable distractor set. The remaining revision failures were three invalid templates. No OpenAI author-call errors or main-study runner errors were recorded, so the low one-shot yield in this panel mainly reflects candidate faults rather than unavailable responses. Model-specific outcomes appear in Appendix~\ref{app:comparisons}, Table~\ref{tab:panel-models}.

One technically passing candidate in each checked strategy was withheld after a difficulty mismatch. Thus, revision produced thirty technical passes but twenty-nine delivered passes. One-shot has no delivery gate, so its delivery rows are inapplicable. The recorded request checks covered parameter restrictions and specified programming features; unrestricted wording requirements and whether a feature was central to the solution still require review.

The higher yield required additional work. Standalone strategy accounting charged the shared initial call and checks to each applicable strategy and excluded readiness probes. Recorded total token usage, including judges, was 244,655 for one-shot, 757,425 for regeneration and 906,530 for revision. Revision used more tokens overall than regeneration, although its tokens per technical pass were lower (about 30,218 versus 34,428). Tokens have different prices across models, and these figures are not monetary savings. Shared judge traffic and the observed service window also affect timings.

\subsection{Controlled repair and hosted-service reliability}\label{sec:repair-results}

A separate experiment started from forty deliberately faulty templates, covering incorrect answer recipes, colliding distractors, correct answers included as distractors and missing requested features. Both strategies started from the same injected candidate and had up to three new candidate opportunities. Revision received the candidate and failure evidence; regeneration produced fresh proposals without that evidence. Difficulty judging was omitted to focus on technical repair. Success required passing the offline checks and preserving the recorded request constraints.

The time-limited run completed 97 of 120 planned model--fixture comparisons. Revision succeeded in 95/97 (97.9\%), compared with 76/97 (78.4\%) for regeneration. Regeneration had seventy-seven technical passes, but one changed the requested parameter domains. Both failed revisions still had distractor faults. Uneven model coverage and deliberately introduced faults limit how far this result can explain ordinary authoring. The twenty-three unrun comparisons are incomplete observations, not failed repairs. Appendix~\ref{app:comparisons} gives the breakdown.

A concurrent SoCLaaS study used the same fifteen requests per configuration but completed only 21/30 comparisons: all fifteen GLM requests and six Qwen requests. Among completed comparisons, one-shot, regeneration and revision produced seven, twelve and thirteen technical passes respectively. Two revision passes changed recorded constraints, leaving eleven constraint-preserving passes. This partial study is reported separately from the complete OpenAI panel.

Its main-study call records contain ten authoring timeouts, nine empty authoring responses, five output-budget failures and one judge timeout. These are call events, not counts of faulty questions. Authoring requested low reasoning and a locally configured 16,384-token allowance; the judge retained provider-default reasoning. Reaching that allowance means the configured response budget was exhausted, not that a SoCLaaS software defect has been identified. The nine unfinished comparisons and shared judge traffic prevent a clean comparison of provider effectiveness or speed.

\subsection{Template reuse}\label{sec:reuse}

The earlier reuse study generated batches of one, five and ten drafts from saved templates. Its upfront costs included unsuccessful jobs before a usable artifact was obtained. For example, the first functions authoring job ended with a difficulty-assessment timeout and the next succeeded.

Saved artifacts generated further drafts without model calls. Variation remained limited by their input domains: the functions template allowed four input combinations, and ten drafts contained only three distinct questions. Separate execution checks assessed the generated answers. Appendix~\ref{app:comparisons}, Table~\ref{tab:reuse} records these demonstrations and their original configuration. Local reuse removes further inference for expansion, but overall cost also depends on authoring failures, duplicates and review effort.

\subsection{Difficulty assessment and authoring policy}\label{sec:difficulty-results}

The v0.3 reference contains thirty Python tracing questions: fifteen easy, nine medium and six hard. The labels were reviewed manually, but all questions informed rubric development and no learner-response data were collected. The following figures measure agreement with that project reference. They are not held-out predictions of student performance.

\subsubsection{Structured judging}\label{sec:structured-difficulty}

The main summary covers four selected hosted configurations, each assessing the thirty questions twice, giving 240 recorded assessments. The reference level was hidden, and the models returned dimension ratings, a trace-event count and an overall label. Table~\ref{tab:difficulty} distinguishes agreement among usable labels from successful assessment coverage. The original study also evaluated Qwen; all five configurations are reported in Appendix~\ref{app:history}, Table~\ref{tab:full-difficulty}.

\begin{table}[!htbp]
\centering\setstretch{1}\small
\setlength{\tabcolsep}{5pt}\renewcommand{\arraystretch}{1.15}
\caption{Structured difficulty judging across the selected four configurations}\label{tab:difficulty}
\begin{tabularx}{\linewidth}{@{}X r@{}}
\toprule
\textbf{Measure} & \textbf{Result}\\
\midrule
Requested and recorded assessments & 240\\
Usable difficulty labels & 210/240 (87.5\%)\\
Labels matching the reference & 200/210 (95.2\%)\\
Usable, matching labels across all requests & 200/240 (83.3\%)\\
Wrong usable labels & 10\\
Missing or invalid labels & 30\\
Easy--hard confusions among usable labels & 0\\
\bottomrule
\end{tabularx}
\end{table}

Eleven assessments exhausted the 16,384-token output budget before returning a valid response. Eighteen returned no structured difficulty value, although some explanations mentioned a level. One explicitly reported being unable to assess. The experimental schema permitted a null label without requiring a reason, so missing values could pass schema validation. The evaluation retained these as unsuccessful assessments rather than extracting labels retrospectively from prose. These failures occurred under provider-default reasoning, separately from the low-reasoning authoring experiment.

The 95.2\% figure applies only to the selected four configurations and their usable labels. Including the fifth configuration gives 244/270 (90.4\%) agreement among usable labels and 244/300 (81.3\%) across all requested assessments. Appendix~\ref{app:history} retains that full comparison. All thirty missing assessments came from the selected four configurations. The additional configuration returned sixty usable labels, with forty-four reference matches.

\subsubsection{Fast decisions and a decision-first cascade}\label{sec:cascade-results}

A second design restricted models to a single three-class decision, with reasoning disabled and scores obtained from first-token preferences for easy, medium and hard. On the initial design comparison, this approach matched 92/120 reference labels, compared with 76/120 for three binary one-versus-rest calls. The designs differ in both decision structure and call count.

The full four-configuration run returned labels for all 240 assessments and matched 184 (76.7\%; balanced accuracy 69.2\%). Its median call time was 493\,ms. Predictions agreed in 116/120 question--configuration pairs across temperatures 0 and 0.7. All fifty-six mistakes were between adjacent levels, including thirty hard questions labelled medium. This makes the middle boundary a substantial limitation despite the absence of easy--hard confusions.

Higher token-preference confidence was associated with greater agreement. A threshold of 0.80 retained 87/240 decisions, of which 75 matched; 0.90 retained 24, all matching. These thresholds were examined on the same reference and remain exploratory. Token preference is a different signal from a model's self-reported confidence.

A replay study then referred low-confidence decisions to the saved structured GLM judge. Below confidence 0.70, the judge matched 84/95 reference labels, compared with 60/95 for the fast decisions. This motivated a live cascade using x-test-3 first and x-test-1 as the reasoning judge, with threshold 0.70. Of thirty questions, sixteen were accepted at the first stage and all matched the reference. Fourteen were escalated; twelve matched and two produced no usable judge label. Overall agreement was therefore 28/30 (93.3\%), counting the unresolved outcomes as unsuccessful.

The cascade used forty-four calls: thirty first-stage calls and fourteen judge calls. Median final latency was 1.50 seconds and mean latency 36.92 seconds, reflecting slow escalated cases. This was a live test on the existing development reference, not an unseen validation set. It supports further integration work rather than a claim of calibrated difficulty prediction.

\subsubsection{Policy used in the authoring comparison}\label{sec:difficulty-policy}

The reported authoring study used the existing v0.2 structured judge, not the decision-first cascade or the v0.3 research contract. A matching estimate passed. A mismatch with self-reported confidence of at least 0.80 requested revision; a lower-confidence mismatch could pass with a warning. Judge errors prevented workflow completion. The v0.3 agreement figures and token-based confidence experiments therefore do not validate the production threshold. Learner studies would be needed to connect either rubric to experienced difficulty.

\subsection{Learning-objective relevance}\label{sec:relevance-results}

The latest comparison rebuilt chemistry question--objective pairs from the SmartSTEM source used by \textcite{zur2023}. Questions were split before evaluation, with 92 development pairs and 36 test pairs. All pairs for a question stayed in one split. Each method's cutoff was fitted on development accuracy and then applied unchanged to the test set. The test set contained twelve relevant and twenty-four non-relevant pairs.

Embedding similarity was compared with the local Qwen3 reranker and four hosted decision configurations. The hosted models were instructed to answer yes or no without reasoning, and the score used their first-token yes/no preferences. Table~\ref{tab:relevance} keeps the methods separate because an aggregate would hide their differences.

\begin{table}[!htbp]
\centering\setstretch{1}\small
\setlength{\tabcolsep}{5pt}\renewcommand{\arraystretch}{1.15}
\caption{Relevance decisions on 36 test pairs}\label{tab:relevance}
\begin{tabularx}{\linewidth}{@{}X r r r@{}}
\toprule
\textbf{Method/configuration} & \textbf{Correct} & \textbf{False accept} & \textbf{False reject}\\
\midrule
Embedding similarity & 32/36 (88.9\%) & 1 & 3\\
Local Qwen3 reranker, 0.6B & 34/36 (94.4\%) & 1 & 1\\
Decision: Qwen3.8 27B & 35/36 (97.2\%) & 1 & 0\\
Decision: Qwen3.8 Flash Next & 35/36 (97.2\%) & 1 & 0\\
Decision: DeepSeek alias & 34/36 (94.4\%) & 1 & 1\\
Decision: Ornith alias & 32/36 (88.9\%) & 3 & 1\\
\bottomrule
\end{tabularx}
\end{table}

The two Qwen decision configurations made one fewer error than the local reranker and three fewer than embeddings. The test is small: those differences support further comparison, rather than a general ranking of relevance methods. The cutoff is fitted to this dataset, and neither cosine similarity nor normalised yes/no preference is a calibrated probability of educational alignment.

The source assigns textbook-subsection objectives to questions, so positive pairs are not individually verified statements that a skill is necessary to solve a question. Negative pairs use other objectives, often from the same chapter. A question-level split avoids reusing the same question across development and test, but it does not create an independent objective taxonomy. The pilot measures agreement on those chemistry pairs; Python learning-objective alignment has not yet been evaluated.

The earlier pilot and adopted local configuration used different instructions and cutoffs. Their 47/54 and 28/36 results remain separate in Appendix~\ref{app:history}. The latest 34/36 reranker result belongs to the rebuilt-pair protocol, not the deployed configuration. Relevance was not a required check in the authoring comparison.

\subsection{Checker verification}\label{sec:checker-results}

The code-domain checker accepted fifteen known-good fixtures and detected all fifteen injected wrong-answer recipes and fifteen injected colliding distractor pools. A wide-domain arithmetic fixture allowed 32,768 input combinations but used thirty-six sampled executions during checking. Thirty later drafts used previously unchecked inputs, and separate execution checks agreed with all thirty answer recipes. This verifies selected implementation behaviours; the rate of detecting naturally occurring model faults and correctness over the full domain remain separate questions.
